%% kapa: exam=1 round=32 year=2021 period=1 form=A subject=real_estate subject_ko=부동산학원론 seq=15 number=95
%% source: Q-Net 감정평가사 기출문제 (cst00309) artlSeq=5208397 file=extracted/1008/5208397/A형/2021년도 제32회 감정평가사 1차 1교시 A형.pdf
%% conversion: kordoc 4.24.51 -> tools/gen1.py (auto); review=proofread-sonnet; flags=none
\begin{question}{95}
비율분석법을 이용하여 산출한 것으로 옳지 않은 것은? (단, 주어진 조건에 한하며, 연간 기준임)
\begin{framedtext}
○ 주택담보대출액: 2억원 \par
○ 주택담보대출의 연간 원리금상환액: 1천만원 \par
○ 부동산가치: 4억원 \par
○ 차입자의 연소득: 5천만원 \par
○ 가능총소득: 4천만원 \par
○ 공실손실상당액 및 대손충당금: 가능총소득의 25\% \par
○ 영업경비: 가능총소득의 50\% \par
\end{framedtext}
\begin{choices}
  \item 부채감당률(DCR)= 1.0
  \item 채무불이행률(DR)= 1.0
  \item 총부채상환비율(DTI)= 0.2
  \item 부채비율(debt ratio)= 1.0
  \item 영업경비비율(OER, 유효총소득 기준)= 0.8
\end{choices}
\end{question}
\answer{5}
